% Unidad U014. Biografía estructural completa del carácter de autoescala.

\label{gfp:b:chap:biografia-completa-phi}

\section{Sélection finie du germe d’auto-échelle}

Soit \(\mathcal U_6\subset\{0,\ldots,999\}^6\) le catalogue des
\(468\) émissions distinctes construit dans le
\cref{gfp:b:chap:semillas-emisor-54}, et soit
\[
 \Psi:\mathcal U_6\longrightarrow\mathbb F_3^6,
 \qquad
 \Psi(U)=((-U_j)\bmod3)_{j=1}^{6}.
\]
Pour \(w\in\Psi(\mathcal U_6)\), rappelons les invariants
\[
 n(w)=\#\Psi^{-1}(w),
 \qquad
 M(w)=\sum_{U\in\Psi^{-1}(w)}\operatorname{mult}(U),
\]
et le vecteur d’occupation
\[
 \nu(w)=
 \bigl(\#\{j:w_j=0\},\#\{j:w_j=1\},\#\{j:w_j=2\}\bigr).
\]

\begin{theorem}[Sélection orbitale de l’auto-échelle]
\label{gfp:b:thm:seleccion-orbital-phi}
L’action diédrale sur \(\Psi(\mathcal U_6)\) contient une unique orbite
qui satisfait le prédicat d’auto-échelle
\[
 n(w)=1,
 \qquad
 M(w)=144,
 \qquad
 \nu(w)=(2,2,2).
\]
La jauge interne
\[
 \operatorname{diag}_c=6,
 \qquad
 \operatorname{card}=ES
\]
sélectionne dans cette orbite le mot
\[
 \boxed{w_6^{\varphi}=121200.}
\]
Son unique émission est
\[
 \boxed{U_{\varphi}=(830,943,830,109,252,252),}
\]
de multiplicité \(144\), de signature multiplicative \(555933\), d’orientation
additive \(ES\), de support multiplicatif de cardinal six et de type résiduel
\(A\).
\end{theorem}

\begin{proof}
La projection directe donne
\[
 U_{\varphi}\bmod3=(2,1,2,1,0,0),
 \qquad
 -U_{\varphi}\bmod3=(1,2,1,2,0,0).
\]
L’énumération est effectuée sur les \(243\) mots visibles et leurs
\(43\) orbites, non sur une sélection d’exemples. Le prédicat conjoint
\(n=1\), \(M=144\), \(\nu=(2,2,2)\) laisse une orbite. En son sein, la
diagonale et l’orientation déclarées laissent une émission. À cette étape,
ni l’équation \(x^2-x-1=0\), ni une table décimale, ni le nom
conventionnel de la constante n’ont été introduits.
\end{proof}

\section{Relèvement par blocs et \(1^{\mathrm a}\) frontière coinductive}

Soit \(b_1=w_6^{\varphi}\), et soit
\(\varepsilon=(0,0,1,1)\). Avec les matrices \(L_0,L_1\) construites dans le
\cref{gfp:b:chap:orbitas-elevacion-r36}, nous définissons
\[
 b_{k+1}=b_kL_{\varepsilon_k}\pmod3
 \qquad(1\leq k\leq4),
\]
et concaténons
\[
 w_{6(k+1)}^{\varphi}=w_{6k}^{\varphi}\mathbin{\|}b_{k+1}.
\]
Le calcul matriciel produit, bloc par bloc,
\[
\begin{aligned}
 b_1&=121200,\\
 b_2&=112202,\\
 b_3&=121020,\\
 b_4&=010210,\\
 b_5&=010200,
\end{aligned}
\]
et, par conséquent,
\begin{equation}
 \boxed{
 w_{30}^{\varphi}
 =121200|112202|121020|010210|010200.}
 \label{gfp:b:eq:w30-phi}
\end{equation}
La neutralité duale de la frontière ajoute
\[
 u_5^{\varphi}=102011,
\]
de sorte que
\begin{equation}
 \boxed{
 w_{36}^{\varphi}
 =121200|112202|121020|010210|010200|102011.}
 \label{gfp:b:eq:w36-phi}
\end{equation}

\begin{remark}[Deux usages distincts d’un mot visible]
Le mot \(102011\) est également visible dans la biographie de \(e\). L’égalité
d’une projection n’identifie pas les états enrichis : le canal,
le préfixe, le quotient, la retenue, la phase et la mémoire appartiennent au type
de l’objet. Cette séparation est obligatoire avant de construire tout
caractère réel.
\end{remark}

\section{Du calendrier tritique au multigraphe d’incidence}

\subsection{Règle modale primitive}

Le calendrier tritique primitif est
\[
 (+1,-1,0).
\]
Nous fixons deux modes visibles, \(\Sigma\) et \(\Pi\), et la règle
\begin{equation}
 +1:m\longmapsto\Sigma,
 \qquad
 -1:m\longmapsto\Pi,
 \qquad
 0:m\longmapsto m.
 \label{gfp:b:eq:regla-modal-primitiva-phi}
\end{equation}
Avec la condition de clôture cyclique, l’orbite modale produite par
\eqref{gfp:b:eq:regla-modal-primitiva-phi} est
\[
 \boxed{(\Sigma,\Pi,\Pi).}
\]
Les trois paires consécutives sont
\[
 \Sigma\longrightarrow\Pi,
 \qquad
 \Pi\longrightarrow\Pi,
 \qquad
 \Pi\longrightarrow\Sigma.
\]

\begin{definition}[Matrice d’incidence chronologique]
Pour \(i,j\in\{\Sigma,\Pi\}\), soit
\[
 N_3(i,j)=
 \#\{r\in\mathbb Z/3\mathbb Z:(m_r,m_{r+1})=(i,j)\}.
\]
L’indice \(3\) compte les instants du calendrier primitif ; il ne désigne pas une
puissance matricielle.
\end{definition}

\begin{theorem}[Descente nécessaire de l’incidence modale]
\label{gfp:b:thm:descenso-incidencia-modal-phi}
Dans l’ordre \((\Sigma,\Pi)\),
\begin{equation}
 \boxed{
 N_3=
 \begin{pmatrix}
 0&1\\
 1&1
 \end{pmatrix}.}
 \label{gfp:b:eq:N3-phi}
\end{equation}
Par répétition exacte du calendrier,
\begin{equation}
 \boxed{
 N_9=3N_3,
 \qquad
 N_{27}=9N_3,
 \qquad
 N_{108}=36N_3.}
 \label{gfp:b:eq:incidencias-9-27-108-phi}
\end{equation}
Ces trois égalités sont des égalités de dénombrements d’arêtes ; elles n’affirment
ni \(N_9=N_3^3\), ni \(N_{27}=N_3^9\), ni \(N_{108}=N_3^{36}\).
\end{theorem}

\begin{proof}
\(\Sigma\to\Sigma\) n’apparaît pas ; chacune des trois autres paires apparaît
une fois. Cela détermine les quatre entrées de \(N_3\). Les calendriers de
longueurs \(9\), \(27\) et \(108\) contiennent respectivement \(3\), \(9\)
et \(36\) copies du bloc primitif. Chaque répétition multiplie tous les
dénombrements par le même entier, ce qui prouve
\eqref{gfp:b:eq:incidencias-9-27-108-phi}.
\end{proof}

\subsection{Publication aréale--volumétrique}

L’étiquette modale et le feuillet géométrique sont des objets distincts. La carte
de publication est
\[
 \mathfrak g_{\rm av}(\Sigma)=\mathscr H_{\rm ar},
 \qquad
 \mathfrak g_{\rm av}(\Pi)=\mathscr H_{\rm vol}.
\]
En transportant \eqref{gfp:b:eq:N3-phi} par \(\mathfrak g_{\rm av}\), nous obtenons
l’opérateur d’incidence
\begin{equation}
 \boxed{
 F_{\rm av}=
 \begin{pmatrix}
 0&1\\
 1&1
 \end{pmatrix}.}
 \label{gfp:b:eq:Fav-derivada-phi}
\end{equation}
Ainsi, les deux transitions croisées, l’unique autorétention volumétrique et
l’absence d’autorétention aréale découlent du calendrier ; ce ne sont pas quatre
entrées indépendantes.

\begin{remark}[Portée de la descente]
L’oubli de la phase élémentaire ne produit pas, à lui seul, un quotient de Markov
fort du Topological Phase Kernel, parce que le mode visible ne détermine pas
lequel des trois opérateurs internes agira au pas suivant. Ce qui
descend en revanche sans hypothèse supplémentaire est le multigraphe des arêtes d’un
bloc complet. Son enveloppe minimale de mémoire un a pour matrice
d’adjacence \(F_{\rm av}\).
\end{remark}

\section{Demi-droite positive invariante et unicité de l’auto-échelle}

\begin{theorem}[Demi-droite de Perron dérivée]
\label{gfp:b:thm:rayo-perron-phi}
Le cône positif de \(F_{\rm av}\) contient une unique demi-droite propre. Si nous la
normalisons comme \(v=(1,x)^{\mathsf T}\), alors
\[
 x^2-x-1=0,
 \qquad
 x>0,
\]
et, par conséquent,
\begin{equation}
 \boxed{x=\frac{1+\sqrt5}{2}=\varphi.}
 \label{gfp:b:eq:phi-perron}
\end{equation}
\end{theorem}

\begin{proof}
L’équation \(F_{\rm av}v=\lambda v\) donne, coordonnée par coordonnée,
\[
 x=\lambda,
 \qquad
 1+x=\lambda x.
\]
En éliminant \(\lambda\), on obtient \(x^2-x-1=0\). Ses racines sont
\((1\pm\sqrt5)/2\), et seule la positive appartient à l’intérieur du cône. La
matrice \(F_{\rm av}\) est primitive, puisque \(F_{\rm av}^2\) a toutes ses
entrées positives ; le théorème de Perron--Frobenius garantit que cette demi-droite
positive est simple et unique. Le nom \(\varphi\) est attribué après cette
construction.
\end{proof}

\begin{corollary}[Équation scalaire d’auto-échelle]
\[
 \varphi=1+\frac1\varphi,
 \qquad
 \varphi^2=\varphi+1,
 \qquad
 \varphi^{-1}=\varphi-1.
\]
Ces identités sont des conséquences de \eqref{gfp:b:eq:phi-perron} ; elles ne sélectionnent
ni le germe ni la matrice.
\end{corollary}

\section{Publication de l’auto-échelle \(\varphi\) par les puissances de Fibonacci et la convergence diophantienne}

Soit \(F_0=0\), \(F_1=1\) et
\[
 F_{n+1}=F_n+F_{n-1}\qquad(n\geq1).
\]

\begin{theorem}[Puissances exactes de l’incidence]
\label{gfp:b:thm:potencias-fibonacci-phi}
Pour tout \(n\geq1\),
\begin{equation}
 \boxed{
 F_{\rm av}^{\,n}=
 \begin{pmatrix}
 F_{n-1}&F_n\\
 F_n&F_{n+1}
 \end{pmatrix}.}
 \label{gfp:b:eq:potencias-Fav-fibonacci}
\end{equation}
\end{theorem}

\begin{proof}
Pour \(n=1\), l’identité est la définition de \(F_{\rm av}\). Si elle vaut
pour \(n\), la multiplication à droite par \(F_{\rm av}\) donne
\[
 \begin{pmatrix}
 F_n&F_{n-1}+F_n\\
 F_{n+1}&F_n+F_{n+1}
 \end{pmatrix}
 =
 \begin{pmatrix}
 F_n&F_{n+1}\\
 F_{n+1}&F_{n+2}
 \end{pmatrix},
\]
ce qui est le cas \(n+1\).
\end{proof}

\begin{theorem}[Identité de Cassini]
\label{gfp:b:thm:cassini-phi}
Pour tout \(n\geq1\),
\begin{equation}
 \boxed{F_{n+1}F_{n-1}-F_n^2=(-1)^n.}
 \label{gfp:b:eq:cassini-phi}
\end{equation}
\end{theorem}

\begin{proof}
Le déterminant de \(F_{\rm av}\) est \(-1\). En prenant les déterminants dans
\eqref{gfp:b:eq:potencias-Fav-fibonacci},
\[
 F_{n-1}F_{n+1}-F_n^2
 =\det(F_{\rm av}^{\,n})=(-1)^n.
\]
\end{proof}

Nous définissons \(r_n=F_{n+1}/F_n\) pour \(n\geq1\). De Cassini découle
\[
 r_n-r_{n-1}
 =\frac{F_{n+1}F_{n-1}-F_n^2}{F_nF_{n-1}}
 =\frac{(-1)^n}{F_nF_{n-1}}.
\]
Par conséquent, les approximations alternent autour de la demi-droite de Perron.

\begin{proposition}[Encadrements rationnels de Perron]
\label{gfp:b:prop:horquillas-perron-phi}
Pour \(m\geq1\),
\begin{equation}
 \frac{F_{2m+2}}{F_{2m+1}}
 <\varphi<
 \frac{F_{2m+1}}{F_{2m}},
 \label{gfp:b:eq:horquilla-fibonacci-phi}
\end{equation}
et la largeur est
\begin{equation}
 \boxed{
 \frac{F_{2m+1}}{F_{2m}}
 -\frac{F_{2m+2}}{F_{2m+1}}
 =\frac{1}{F_{2m}F_{2m+1}}.}
 \label{gfp:b:eq:anchura-fibonacci-phi}
\end{equation}
\end{proposition}

\begin{proof}
L’alternance précédente place les quotients pairs en dessous et les impairs
au-dessus de \(\varphi\). La formule de la largeur est
\eqref{gfp:b:eq:cassini-phi} avec \(n=2m+1\), après réduction des deux fractions au
même dénominateur. Comme \(F_n\to\infty\), la largeur tend vers zéro.
\end{proof}

\section{Incidence de mémoire d’ordre \(1\) et mesure de Parry de l’enveloppe symbolique}

Soit \(X_{\rm av}\) le plus petit sous-décalage de type fini sur
\(\{\mathscr H_{\rm ar},\mathscr H_{\rm vol}\}\) qui admet exactement
les trois paires observées dans le calendrier modal. Sa matrice d’adjacence est
\(F_{\rm av}\).

\begin{theorem}[Transition et mesure de Parry]
\label{gfp:b:thm:parry-phi}
Soit \(r=(1,\varphi)^{\mathsf T}\). La matrice stochastique de Parry est
\begin{equation}
 \boxed{
 P_{ij}=
 \frac{(F_{\rm av})_{ij}r_j}{\varphi r_i}
 =
 \begin{pmatrix}
 0&1\\
 \varphi^{-2}&\varphi^{-1}
 \end{pmatrix}.}
 \label{gfp:b:eq:parry-matriz-phi}
\end{equation}
La mesure stationnaire des sommets est proportionnelle à
\((1,\varphi^2)\), donc
\begin{equation}
 \boxed{
 \frac{\mu_{\rm ar}}{\mu_{\rm vol}}=\varphi^{-2}.}
 \label{gfp:b:eq:parry-razon-phi}
\end{equation}
\end{theorem}

\begin{proof}
La somme de la ligne \(i\) de \(P\) est
\[
 \frac{(F_{\rm av}r)_i}{\varphi r_i}=1.
\]
Comme \(F_{\rm av}\) est symétrique, les vecteurs propres gauche et droit
coïncident. Les poids stationnaires sont proportionnels à leurs produits
coordonnée par coordonnée, c’est-à-dire à \((1,\varphi^2)\). Le rapport entre les
deux poids est \(\varphi^{-2}\).
\end{proof}

\begin{remark}[Trois mesures à ne pas confondre]
La fréquence chronologique de l’orbite primitive est
\[
 \nu_{\rm cyc}(\mathscr H_{\rm ar})=\frac13,
 \qquad
 \nu_{\rm cyc}(\mathscr H_{\rm vol})=\frac23.
\]
La mesure uniforme \(\tau_{468}\) vit sur le catalogue des émissions. La
mesure de Parry vit sur les histoires admissibles de longueur arbitraire. La
première compte les instants, la deuxième compte les émissions et la troisième pondère
les cylindres dynamiques. Aucune n’est une approximation des autres.
\end{remark}

\section{Caractère multiplicatif et conservation cylindrique}

Pour une route admissible \(\gamma:i\to j\) de longueur \(n\), nous définissons
\begin{equation}
 \chi_P(\gamma)=\varphi^n\frac{r_i}{r_j}.
 \label{gfp:b:eq:caracter-parry-phi}
\end{equation}

\begin{proposition}[Multiplicativité par concaténation]
\label{gfp:b:prop:multiplicatividad-parry-phi}
Si \(\gamma_1:i\to j\) et \(\gamma_2:j\to k\), alors
\[
 \chi_P(\gamma_2\circ\gamma_1)
 =\chi_P(\gamma_2)\chi_P(\gamma_1).
\]
Sur une route fermée de longueur \(n\),
\[
 \chi_P(\gamma)=\varphi^n.
\]
\end{proposition}

\begin{proof}
La longueur s’additionne et la coordonnée intermédiaire se simplifie par télescopage :
\[
 \varphi^{n_1+n_2}\frac{r_i}{r_k}
 =\left(\varphi^{n_1}\frac{r_i}{r_j}\right)
  \left(\varphi^{n_2}\frac{r_j}{r_k}\right).
\]
Si \(i=k\), le quotient terminal est égal à un.
\end{proof}

Nous normalisons le vecteur gauche comme
\[
 \ell=\frac{r}{1+\varphi^2},
 \qquad
 \ell^{\mathsf T}r=1.
\]
La mesure d’un cylindre est
\[
 \mu[x_0\cdots x_n]
 =\ell_{x_0}r_{x_n}\varphi^{-n}.
\]
Par conséquent,
\begin{equation}
 V_n(x)
 :=\frac{\mu[x_0]}{\mu[x_0\cdots x_n]}
 =\varphi^n\frac{r_{x_0}}{r_{x_n}}
 =\chi_P(x_0\cdots x_n).
 \label{gfp:b:eq:volumen-cilindrico-phi}
\end{equation}
Pour une dimension typée \(D>0\), l’échelle
\[
 a_n^{(D)}(x)=V_n(x)^{1/D}
\]
satisfait la loi exacte
\begin{equation}
 \boxed{
 \bigl(a_n^{(D)}(x)\bigr)^D
 \mu[x_0\cdots x_n]=\mu[x_0].}
 \label{gfp:b:eq:conservacion-cilindrica-phi}
\end{equation}

En dimension trois et sur des retours fermés,
\[
 \frac{a_9}{a_0}=\varphi^3,
 \qquad
 \frac{a_{27}}{a_0}=\varphi^9,
 \qquad
 \frac{a_{108}}{a_0}=\varphi^{36}.
\]
Si \(A_k=a_{9k}\), alors
\begin{equation}
 A_{k+1}-2A_k+A_{k-1}
 =A_{k-1}(\varphi^3-1)^2>0.
 \label{gfp:b:eq:convexidad-retornos-phi}
\end{equation}
C’est une convexité en profondeur de retour ; elle n’est pas interprétée comme une
grandeur cinématique sans une carte supplémentaire de durée.

\section{Mémoire quadratique et branche contractée}

Les valeurs propres de \(F_{\rm av}\) sont
\[
 \varphi,
 \qquad
 -\varphi^{-1}.
\]
La branche contractée change d’orientation à chaque itération. Pour conserver
ce signe avant d’élever au carré, nous passons à
\(\operatorname{Sym}^2(\mathbb R^2)\). Dans la base
\((x^2,2xy,y^2)\),
\begin{equation}
 \operatorname{Sym}^2(F_{\rm av})=
 \begin{pmatrix}
 0&0&1\\
 0&1&2\\
 1&1&1
 \end{pmatrix}.
 \label{gfp:b:eq:sym2-Fav-phi}
\end{equation}

\begin{theorem}[Spectre de la mémoire de second ordre]
\label{gfp:b:thm:espectro-sym2-phi}
\[
 \det\bigl(tI-\operatorname{Sym}^2(F_{\rm av})\bigr)
 =(t+1)(t^2-3t+1),
\]
et
\begin{equation}
 \boxed{
 \operatorname{Spec}\bigl(\operatorname{Sym}^2(F_{\rm av})\bigr)
 =\{\varphi^2,-1,\varphi^{-2}\}.}
 \label{gfp:b:eq:espectro-sym2-phi}
\end{equation}
L’unique mode stable positif est \(\varphi^{-2}\).
\end{theorem}

\begin{proof}
Les valeurs propres d’une puissance symétrique de degré deux sont les produits
\(\lambda_1^2\), \(\lambda_1\lambda_2\) et \(\lambda_2^2\). Avec
\(\lambda_1=\varphi\) et \(\lambda_2=-\varphi^{-1}\), on obtient
\(\varphi^2\), \(-1\) et \(\varphi^{-2}\). La première dépasse un, la
deuxième a pour module un et la troisième appartient à \((0,1)\).
\end{proof}

Sur une route fermée avec \(V_n=\varphi^n\), le mode stable satisfait
\begin{equation}
 \boxed{M_n=\varphi^{-2n}=V_n^{-2}.}
 \label{gfp:b:eq:modo-estable-phi}
\end{equation}

\section{Réalisation hyperbolique normalisée de l’auto-échelle : \(\exp(2\log\varphi\,J_h)=F_{\mathrm{av}}^2\)}

Soit
\[
 A=F_{\rm av}^{\,2}=
 \begin{pmatrix}1&1\\1&2\end{pmatrix},
 \qquad
 J_h=\frac{2A-3I}{\sqrt5}.
\]
Alors
\[
 J_h^2=I,
 \qquad
 A=\frac32I+\frac{\sqrt5}{2}J_h.
\]
Comme
\[
 \cosh(2\log\varphi)=\frac32,
 \qquad
 \sinh(2\log\varphi)=\frac{\sqrt5}{2},
\]
on obtient
\begin{equation}
 \boxed{
 \exp(2\log\varphi\,J_h)=F_{\rm av}^{\,2}.}
 \label{gfp:b:eq:realizacion-hiperbolica-phi}
\end{equation}
La constante d’auto-échelle est donc simultanément la valeur propre de
Perron de l’incidence et le paramètre exponentiel de sa réalisation
hyperbolique normalisée.

\section{Caractérisation, encadrements et prolongement}

La construction finie détermine la demi-droite positive de \(F_{\rm av}\). La
caractérisation analytique ultérieure est
\[
 x^2-x-1=0,
 \qquad x>0.
\]
Les encadrements \eqref{gfp:b:eq:horquilla-fibonacci-phi} ont des extrémités rationnelles
et une largeur explicite \eqref{gfp:b:eq:anchura-fibonacci-phi}. Pour chaque niveau
nonadique \(K\), nous choisissons le plus petit \(m=m_{\varphi}(K)\) dont l’encadrement est
contenu dans un unique sous-cylindre parmi les \(729\) éléments de la fibre
ambiante. La minimalité fait de la localisation une fonction, non un
choix tacite.

Les cylindres sélectionnés sont compatibles avec les troncatures et leurs
diamètres tendent vers zéro. Leur intersection contient un unique réel ; le
\cref{gfp:b:thm:rayo-perron-phi} l’identifie à \(\varphi\). Les quotients de
Fibonacci sont donc des certificats rationnels de la localisation, non des données
décimales employées pour décider du germe.

\Needspace{22\baselineskip}
\section{Certificat interne d’unicité, de compatibilité et de convergence de l’auto-échelle de Perron}

\begin{remark}[Certificat et contrôle de validité : réfutation de la biographie d’auto-échelle]
La construction est réfutée si l’un quelconque des faits
suivants se produit :
\begin{enumerate}
 \setlength{\itemsep}{0.15\baselineskip}
 \item le prédicat fini sélectionne plus d’une orbite ou plus d’une émission ;
 \item la récurrence \(L_0,L_0,L_1,L_1\) ne produit pas
       \eqref{gfp:b:eq:w30-phi} ;
 \item le calendrier ne produit pas exactement les trois paires utilisées dans
       \eqref{gfp:b:eq:N3-phi} ;
 \item un multiple de dénombrement \(N_{9},N_{27},N_{108}\) est confondu avec une
       puissance de \(N_3\) ;
 \item \(F_{\rm av}\) possède une autre demi-droite propre positive ;
 \item \eqref{gfp:b:eq:potencias-Fav-fibonacci} ou l’identité de Cassini échoue ;
 \item les encadrements rationnels ne sont pas emboîtés ou ne contractent pas leur largeur ;
 \item la fréquence chronologique, la mesure uniforme du
       catalogue ou la mesure de Parry sont identifiées ;
 \item une table décimale intervient avant la construction de la demi-droite de Perron ;
 \item un cylindre sélectionné ne se tronque pas au cylindre précédent.
\end{enumerate}
\end{remark}
